\documentclass[10pt]{article}
\usepackage{titling}

\title{Ethers}
\date{DATE PLACEHOLDER}
\author{Ingyu Bahng}

\usepackage[letterpaper, margin=1in]{geometry}
\usepackage{amsmath} % better math functions
\usepackage{amssymb} % better math symbols
\usepackage{babel}[english] % change rules to english
\usepackage{lipsum} % allows for lorem ipsum text generation
\usepackage{xr-hyper} % allows cross referencing labels from external LaTeX documents 
\usepackage{hyperref} % for hyperlinking text
\usepackage{fancyvrb} % for better verbatim environments
\usepackage{newverbs} % allows for inline typewriter font via \texttt{}
\usepackage{fancyvrb} % also code font blocks but has more capabilities than texttt
\usepackage{footnote} % allows footnotes inside tabulars
\usepackage{dcolumn} % allows for decimal alignment in tables
\usepackage{tabularx}

\usepackage{fontspec}
\setmainfont{Gentium Plus} % allowing greek letters in text

\usepackage{unicode-math}
\setmathfont{XITS Math} % allowing greek letters in math



% tcolorbox package
\usepackage{tcolorbox}
\tcbuselibrary{breakable}

\tcbset{ % this is the default design of tcolorboxes
  colframe = black,
  colback  = black!1,
  coltitle = black,
  colbacktitle = black!15,
  breakable, 
  arc=0mm,
  boxrule=0.5pt,
  toptitle=3pt,
  bottomtitle=3pt,
  left=8pt,
  right=8pt,
  top=6pt,
  bottom=6pt,
  before skip=12pt,
  after skip=12pt,
  bottomrule at break=-1pt,
  toprule at break=-1pt,
  fonttitle=\bfseries,
}

% Custom Environments
\newtcolorbox[auto counter, number within=section]{definition}[1][]
{
    title = \textbf{Definition \thetcbcounter: ~#1}
}

\newtcolorbox[auto counter, number within=section]{core}[1][]
{
    title = \textbf{Core Concept \thetcbcounter: ~#1}
}


\usepackage{fancyhdr}
\usepackage[skip=0.8em, indent=0em]{parskip}

\pagestyle{fancy} % this section generates the lines and text at the top and bottom of each page
\fancyhead[L]{\thetitle}
\fancyhead[C]{\theauthor}
\fancyhead[R]{\thedate}
\fancyfoot[C]{\thepage}
\renewcommand{\footrulewidth}{0.3pt}
\renewcommand{\thispagestyle}[1]{}

\renewcommand{\arraystretch}{1.5} % table vertical padding
\setlength{\tabcolsep}{10pt} % table horizontal padding

\usepackage{enumitem} % better control over enumerate and itemize

% List customization
\setlist[itemize]{%
  label=\bullet,
  itemsep=2pt,
  leftmargin=1.5em,
}

\setlist[enumerate]{%
  label=(\alph*),
  itemsep=2pt,
  leftmargin=1.5em,
}



\usepackage{pgfplots} % for plot creation
\pgfplotsset{compat=1.18} % setting the version used for pgfplots
\usepackage{float} % for better figure placement using [h]
\usepackage{graphicx} % for importing figures into the tex file
\usepackage{subcaption} % allows for multiple subfigures with individual captions wihtin one figure
\usepackage{xparse} % allows for more than 9 parameters in commands
\usepackage{adjustbox} % for valign option
\captionsetup{font=small} % setting the caption fontsize to small always

\usepackage{silence}
\WarningsOff[float]  % suppress float package warnings

\newcommand{\myfig}[4]{%
  \begin{figure}[H]
    \centering
    \adjustbox{valign=c}{\includegraphics[width=#1\textwidth]{#2}}
    \caption{#3}
    \label{#4}
  \end{figure}
}

\newcommand{\mymultifig}[8]{
  \begin{figure}[H]%
    \centering%
    \begin{subfigure}{#1\textwidth}%
      \centering%
      \includegraphics[width=\linewidth]{#2}%
    \end{subfigure}%
    \hfill%
    \begin{subfigure}{#3\textwidth}%
      \centering%
      \includegraphics[width=\linewidth]{#4}%
    \end{subfigure}%
    \hfill%
    \begin{subfigure}{#5\textwidth}%
      \centering%
      \includegraphics[width=\linewidth]{#6}%
    \end{subfigure}%
    \caption{#7}%
    \label{#8}%
  \end{figure}%
}

\usepackage{newverbs} % allows for inline typewriter font via \texttt{}
\usepackage{fancyvrb} % also code font blocks but has more capabilities than texttt
\usepackage{listings} % allows for codeblocks - config below

%BELOW ARE CODE BLOCK DESIGNS (Light Theme)
\definecolor{gray}{HTML}{707070}
\definecolor{lightgray}{HTML}{333333}
\definecolor{codebg}{HTML}{F5F5F5}
\definecolor{keyword}{HTML}{0000FF}
\definecolor{string}{HTML}{A31515}
\definecolor{number}{HTML}{098658}
\definecolor{function}{HTML}{795E26}
\definecolor{comment}{HTML}{008000}
\definecolor{classname}{HTML}{267F99}

% default options for listings (for code)
\lstset{
  frame=ltbr,
  language=python,
  aboveskip=3mm,
  belowskip=3mm,
  showstringspaces=false,
  columns=fullflexible,
  keepspaces=true,
  basicstyle=\fontsize{9}{10}\ttfamily\color{lightgray},
  numbers=left,
  firstnumber=1,
  numberstyle=\fontsize{7}{10}\ttfamily\color{gray},
  keywordstyle=\color{keyword},
  commentstyle=\color{comment},
  stringstyle=\color{string},
  backgroundcolor=\color{codebg},
  breaklines=true,
  breakatwhitespace=true,
  tabsize=2,
  xleftmargin=3em,
  numbersep=1em,
  framexleftmargin=2.5em,
  framesep=6pt,
  stepnumber=1,
}

\hfuzz=5.002pt % ignore overfull hbox badness warnings below this limit



\usepackage{chemfig}
\usepackage[version=4]{mhchem}

\newcommand{\bce}[1]{
  \hspace{7mm}{#1}
}

\newcommand{\sno}[0]{
  $\mathrm{S}_{\mathrm{N}}1$
}

\newcommand{\snt}[0]{
  $\mathrm{S}_{\mathrm{N}}2$
}

\newcommand{\reactionfig}[4]{
    \begin{figure}[H]
    \centering
    \scalebox{#4}{
    \schemestart
    #1
    \schemestop
    }
    \caption{#2}
    \label{#3}
    \end{figure}
}




\externaldocument[gas-]{../gases/gases}[../gases/gases.pdf]
\externaldocument[aci-]{../acids_and_bases/acids_and_bases}[../acids_and_bases/acids_and_bases.pdf]
\externaldocument[equ-]{../equilibrium/equilibrium}[../equilibrium/equilibrium.pdf]
\externaldocument[the-]{../thermodynamics/thermodynamics}[../thermodynamics/thermodynamics.pdf]
\externaldocument[qua-]{../quantization/quantization}[../quantization/quantization.pdf]
\externaldocument[alc-]{../alcohols/alcohols}[../alcohols/alcohols.pdf]
\externaldocument[nom-]{../nomenclature/nomenclature}[../nomenclature/nomenclature.pdf]
\externaldocument[rad-]{../radical_halogenation/radical_halogenation}[../radical_halogenation/radical_halogenation.pdf]
\externaldocument[con-]{../conformational_analysis/conformational_analysis}[../conformational_analysis/conformational_analysis.pdf]
\externaldocument[alk-]{../alkane_alkene_alkyne/alkane_alkene_alkyne}[../alkane_alkene_alkyne/alkane_alkene_alkyne.pdf]
\externaldocument[nuc-]{../nuc_sub_and_elim/nuc_sub_and_elim}[../nuc_sub_and_elim/nuc_sub_and_elim.pdf]
%\externaldocument[eth-]{../ethers/ethers}[../ethers/ethers.pdf]

\begin{document}

\input{../../presets/_general/startup}

\section{Ether Synthesis}

  \begin{definition}{Ethers}
 
    Ethers are a class of organic compounds characterized by an oxygen bonded to two alkyl or aryl groups. Their general formula is R—O—R', where $R = R'$ (symmetrical ether) or $R \neq R'$ (asymmetrical ether).

    \reactionfig{
      \scheme{
        \chemfig{
          [:0]\ce{^{\textcolor{Blue}{δ^+}}CH3}
          -[:30]\textcolor{Red}{\ce{O^{δ^-}}}
          -[:-30]\ce{CH3^{\textcolor{Blue}{δ^+}}}
        }
      }
    }{Dimethyl ether structure}{filler1}{1.2}

  \end{definition}

  \begin{core}{Williamson Ether Synthesis}
   
    \textbf{Williamson ether synthesis} is the primary way chemists make asymmetrical ethers. This method uses an alkoxide \ce{R{'}O-} as its nucleophile, thus works best when attacking non-β-branched primary \ce{R-X}s because β-branched, secondary, or tertiary \ce{R-X}s will unintendedly use the E2 mechanism.

    \bce{\ce{R_{prim}-X + R{'}O- ->[\snt] R-O-R{'}}}

    It's important to note that the if the desired product is asymmetrical, the choice of corresponding R groups of the haloalkane and alkoxdide is important because a tertiary or secondary haloalkane paired with a primary alkoxide will lead to E2 whereas the opposite combination will undergo \snt (\textit{the desired method}).

    \doublereactionfig{
      \scheme{
        \chemfig{
          [:0]
          -[:60]
          (-[:120])
          -[:0]\textcolor{Red}{\ce{O-}}
        }
        \+
        \chemfig{
          [:0]
          -[:30]
          -[:-30]\textcolor{Green}{Br}
        }
        \arrow{->[
          \footnotesize\chemfig{
            [:0]\ce{H3C}
            -[:0]S
            (=[:90]O)
            -[:0]\ce{CH3}
          }
        ][85\%]}
        \chemfig{
          [:0]
          -[:60]
          (-[:120])
          -[:0]\textcolor{Red}{\ce{O}}
          -[:60]
          -[:0]
        }
      }
    }{
      \scheme{
        \chemname{
          \chemfig{
            [:0]
            -[:60]
            (-[:120])
            -[:0]\textcolor{Green}{Br}
          }
        }{\textit{secondary central carbon}}
        \+
        \chemfig{
          [:0]
          -[:30]
          -[:-30]\textcolor{Red}{\ce{O-}}
        }
        \arrow{->}
        \chemfig{
          \text{\textit{E2 Reaction}}
        }
      }
    }{Williamson ether syntehsis retrosynthesis; \textit{unsymmetrical ethers have two retrosynthetic choices, one is usually favorable over the other}}{filler3}{0.9}

  \end{core}
  
  \begin{core}{Alcoholysis of \ce{R_{sec/tert}X}}
    
    For \ce{R_{sec/tert}}X haloalkanes, chemists use the protonated version of alkoxides (alcohol) to make the nucleophile less basic and consequently undergo \sno or E1 rather than E2.

    \reactionfig{
      \scheme{
        \chemfig{
          [:0]\ce{H3C}
          -[:0]C
          (-[:90]\ce{CH3})
          (-[:-90]\ce{CH3})
          -[:0]\textcolor{Green}{\ce{Cl}}
        }
        \+
        \chemfig{
          \textcolor{Red}{\ce{CH3O}}\ce{H}
        }
        \arrow{<=>}
        \chemfig{
          [:0]\ce{H3C}
          -[:0]C
          (-[:90]\ce{CH3})
          (-[:-90]\ce{CH3})
          -[:0]\textcolor{Red}{\ce{OCH3}}
        }
        \+
        \chemfig{
          [:0]\ce{H}\textcolor{Green}{\ce{Cl}}
        }
      }
    }{Alcoholysis of \ce{R_{sec/tert}X}}{fig:filler4}{0.9}

    While E1 reactions still remain an unintended pathway, chemists minimize this by carrying out the reaction in lower temperatures.

  \end{core}

  \begin{core}{Intramolecular Williamson Ether Synthesis (Cyclic Ether Formation)}

    So far, we've looked at only \textbf{intermolecular} ether synthesis reactions. \textbf{Intramolecular} ether synthesis (\textit{or cyclic ether formation} works in the same way as the intermolecular equivalent does.

    \doublereactionfig{
      \scheme{
        \chemfig{
          \textcolor{Blue}{\ce{H}}\textcolor{Red}{\ce{O}}\ce{CH2CH2}\textcolor{Green}{Br}
        }
        \+
        \chemfig{
          \textcolor{Red}{\ce{HO-}}
        }
        \arrow{->}
        \chemfig{
          *3(-
          \textcolor{Red}{O}
          --)
        }
        \+
        \chemfig{
          \textcolor{Green}{\ce{Br-}}
        }
        \+
        \chemfig{
          \textcolor{Blue}{\ce{H}}\textcolor{Red}{\ce{OH}}
        }
      }
    }{
      \scheme{
        \chemfig{
          \textcolor{Blue}{\ce{H}}\textcolor{Red}{\ce{O}}\ce{(CH2)4CH2}\textcolor{Green}{Br}
        }
        \+
        \chemfig{
          \textcolor{Red}{\ce{HO-}}
        }
        \arrow{->}
        \chemfig{
          *6(---
          \textcolor{Red}{O}
          ---)
        }
        \+
        \chemfig{
          \textcolor{Green}{\ce{Br-}}
        }
        \+
        \chemfig{
          \textcolor{Blue}{\ce{H}}\textcolor{Red}{\ce{OH}}
        }
      }
    }{Intramolecular Williamson ether synthesis}{fig:filler5}{0.9}

    \textbf{Intramolecular are particularly favorable over intermolecular ether synthesis} because [1] it does not lower entropy and [2] highly diluted concentrations of the relevant alcohol do not effect it as much because the reaction is occurring with itself.

    Mechanistically, this reaction occurs via a backside attack with inversion, prepared by the rearrangement of a less stable rotamer if needed.

    \reactionfig{
      \scheme{
        \chemname{
          \chemfig{
            [:0]\ce{H3C}
            -[:60]C
            (<:[:100]H)
            (<[:170]Br)
            -[:0]C
            (<[:-80]H)
            (<:[:-10]H)
            -[:60]\ce{OH}
          }
        }{\textit{most stable rotamer, but unsuitable for \snt}}
        \arrow{<=>[\textit{rotation}][]}
        \chemfig{
          [:0]Br
          -[:60]C
          (<:[:100]\ce{H3C})
          (<[:170]H)
          -[:0]C
          (<[:-80]H)
          (<:[:-10]H)
          -[:60]\ce{OH}
        }
        \arrow{->[THF, \ce{K2CO3}][]}
        \chemfig{
          [:0]@{bromine}\textcolor{Green}{Br}
          -[@{bond}:60,,,,Green]@{carbon}C
          (<:[:100]\ce{H3C})
          (<[:170]H)
          -[:0]C
          (<[:-80]H)
          (<:[:-10]H)
          -[:60]@{oxygen}\textcolor{Red}{\ce{O-}}
        }
        \arrow{->[84\%][]}
        \chemfig{
          [:90]*3(
          C
          (<[:-80]H)
          (<:[:-20]H)
          -
          \textcolor{Red}{O}
          -
          C
          (<[:-100]H)
          (<:[:-160]\ce{H3C})
          -)
        }
        \chemmovearrow{
          \reactionfigarrow{oxygen}{180}{carbon}{30}{Red} 
          \reactionfigarrow{bond}{0}{bromine}{-30}{Red} 
        }
      }
    }{Backside attack and inversion in cyclic ether formation}{fig:filler6}{0.7}

    The relative rates of ring closure depend on the trade off between ΔH (\textit{ring strain}) and ΔS (\textit{lining up}) in the various transition states. The following is the ranking of ether ring formation speeds by ring size.

    \[
      3 > 5 \geq 6 > 4 \geq 7 > 8
    \]

    Reaction speed does not directly correspond to ring size. Although one might expect a linear relationship, the trend is more nuanced: 3-membered rings react fastest owing to the close proximity of the bonding atoms, 4-membered rings are the slowest, and 6-membered rings are surprisingly fast again due to favorable atomic geometry. In summary, this ordering stems from competing ΔH and ΔS contributions across the transition states of each ring size, rather than any simple geometric trend.

    \doublereactionfig{
      \scheme{
        \chemname{
          \chemfig{
            [:0]@{l1}\textcolor{Green}{L}
            -[@{b1}:30,,,,Green]@{carbon1}
            -[:-30]
            -[:30]@{o1}\textcolor{Red}{O}
          }
        }{\textcolor{Red}{super fast}}
        \qquad\qquad
        \chemname{
          \chemfig{
            [:0]@{l2}\textcolor{Green}{L}
            -[@{b2}:30,,,,Green]@{c2}
            -[:-30]
            -[:30]
            -[:90]@{o2}\textcolor{Red}{O}
          }
        }{\textcolor{Red}{fast}}
      }
      \chemmovearrow{
        \reactionfigarrow{o1}{135}{carbon1}{45}{Red}
        \reactionfigarrow{b1}{120}{l1}{150}{Red}
        \reactionfigarrow{o2}{210}{c2}{45}{Red}
        \reactionfigarrow{b2}{120}{l2}{150}{Red}
      }
    }{
      \scheme{
        \chemname{
          \chemfig{
            [:0]@{l1}\textcolor{Green}{L}
            -[@{b1}:30,,,,Green]@{carbon1}
            -[:-30]
            -[:30]
            -[:90]
            -[:150]@{o1}\textcolor{Red}{O}
          }
        }{\textcolor{Red}{fast}}
        \qquad\qquad
        \chemname{
          \chemfig{
            [:0]@{l2}\textcolor{Green}{L}
            -[@{b2}:30,,,,Green]@{c2}
            -[:-30]
            -[:30]
            -[:90]
            -[:150]
            -[:-150]@{o2}\textcolor{Red}{O}
          }
        }{\textcolor{Red}{fast}}
      }
      \chemmovearrow{
        \reactionfigarrow{o1}{-120}{carbon1}{90}{Red}
        \reactionfigarrow{b1}{120}{l1}{150}{Red}
        \reactionfigarrow{o2}{-90}{c2}{90}{Red}
        \reactionfigarrow{b2}{120}{l2}{150}{Red}
      }
    }{Relative rates of ring closure}{fig:filler7}{1}

  \end{core} 

  \begin{core}{Symmetrical Ether Synthesis via Alcohols}
    
    Ethers can also be formed by \textbf{alcohols and \ce{H+} through \snt and/or \sno} depending on the R groups. Chemists specifically use this technique to form symmetrical ethers because and alcohol mixure would otherwise produce a slurry of different ether product.

    \bce{\ce{2R-OH ->[H2SO4][-H2O] R-O-R}}

    For \ce{R_{prim}-OH}, this mechanism works via \snt due to minimal steric hinderance.

    \reactionfig{
      \scheme{
        \chemfig{
          R
          @{oxygen1}\textcolor{Red}{O}
          H
        }
        \qquad
        \qquad
        \chemfig{
          [:0,,,,Green]@{rgroup}\textcolor{Black}{R}
          -[@{bond}:0]@{oxygen2}\textcolor{Green}{\pcharge{0}{O}}
          (-[:-60]\textcolor{Green}{H})
          -[:60]\textcolor{Green}{H}
        }
      }
      \chemmovearrow{
        \reactionfigarrow{oxygen1}{45}{rgroup}{90}{Red}
        \reactionfigarrow{bond}{-120}{oxygen2}{-100}{Red}
      }
    }{Symmetrical ether formation via \snt}{fig:filler8}{1}

    When dealing with \ce{R_{sec/tert}-OH}, this mechanism shifts to \sno.

    \reactionfig{
      \scheme{
        \chemfig{
          [:0]
          -[:60]
          (-[:120])
          -[:0]\textcolor{Red}{O}H
        }
        \qquad
        \chemfig{
          [:0]
          -[:60]
          (-[:120])
          -[@{bond1}:0,,,,Green]@{oxygen1}\textcolor{Green}{\pcharge{0}{O}}
          (-[:60,,,,Green]\textcolor{Green}{H})
          -[:-60,,,,Green]\textcolor{Green}{H}
        }
        \arrow{->}
        \chemfig{
          [:0]
          -[:60]
          (-[:120])
          -[:0]@{oxygen2}\textcolor{Red}{O}
          -[@{bond2}:-60]H
        }
        \qquad
        \chemfig{
          [:0]
          -[:60]@{carbon1}\pcharge{0}{}
          -[:120]
        }
        \arrow{->}
        \chemfig{
          [:0]
          -[:60]
          (-[:120])
          -[:0]\textcolor{Red}{O}
          -[:0]
          (-[:60])
          -[:-60]
        }
      }
      \chemmovearrow{
        \reactionfigarrow{bond1}{-100}{oxygen1}{-90}{Red}
        \reactionfigarrow{bond2}{-135}{oxygen2}{-150}{Red}
        \reactionfigarrow{oxygen2}{30}{carbon1}{180}{Red}
      }
    }{Symmetrical ether formation via \sno}{fig:filler9}{1}

  \end{core} 

  \begin{core}[label=core:tbu]{Asymmetrical Ether Synthesis (\textit{Tetrabutyl Scenario})}
    
    However, a unique alcohol combination that preferentially forms an asymetrical ether is an \textbf{unhindered \ce{ROH} and tetrabutyl mixture}.

    \reactionfig{
      \scheme{
        \chemfig{
          [:0]\ce{H3C}
          -[:0]
          (-[:90]\ce{CH3})
          (-[:-90]\ce{CH3})
          -[:0]OH
        }
        \+
        \chemfig{
          \textcolor{Green}{\ce{RCH2}}\ce{OH}
        }
        \arrow{->[\ce{H+}][]}
        \chembrackets{
          \chemfig{
            [:0]\ce{H3C}
            -[:0]@{carbon}\pcharge{0}{C}
            (-[:-60]\ce{CH3})
            -[:60]\ce{CH3}
          }
        }
        \arrow{->[
          \chemfig{
            \textcolor{Green}{\ce{RCH2}}@{oxygen}\textcolor{Red}{O}H
          }
        ][-\ce{H+}]}
        \chemfig{
          [:0]\textcolor{Green}{\ce{RCH2}}\ce{O}
          -[:0]
          (-[:90]\ce{CH3})
          (-[:-90]\ce{CH3})
          -[:0]\ce{CH3}
        }
      }
      \chemmovearrow{
        \reactionfigarrow{oxygen}{90}{carbon}{30}{Red}
      }
    }{Asymmetrical ether formation with t-Butyl group}{fig:filler10}{0.8}

    The \ce{ROH} preferentially attacks readily avaliable t-Bu cations (\textit{due to hyperconjugation}) via \sno. And while the t-Bu cation could be attacked by another t-Bu alcohol molecule, steric hinderance makes this interaction less likely to occur compared to the smaller, less sterically hindered \ce{ROH}.

    In the laboratory, t-Bu is used for \textbf{"protection"} of other alcohols by forming these t-Bu ethers.

  \end{core}
  
\section{Ether Reactions}

  Let's go through a collection of different ether reactions.

  \begin{core}{\ce{R_{prim}-O-R_{prim}} Reaction}
    
    \ce{R_{prim}-O-R_{prim}} are quite unreactive, exhibiting stability to bases, \ce{RLi}, \ce{RMgBr}, and dilute aqueous \ce{H+}. However they are reactive under \ce{H+X-} exposure.

    \reactionfig{
      \scheme{
        \chemfig{
          [:0]\ce{H3CH2C}
          -[:0]@{oxygen1}\textcolor{Red}{O}
          -[:0,,,,Red]\textcolor{Red}{\ce{CH2CH3}}
        }
        \arrow{<=>[
          \chemfig{
            @{proton}\ce{H+}
          }
        ][]}
        \chemfig{
          [:0,,,,Green]\ce{H3C}
          -[:0,0.7,,,Black]@{carbon}C
          (-[:-90,0.5,,,Black]H)
          (-[:90,0.5,,,Black]H)
          -[@{bond}:0]@{oxygen2}\textcolor{Green}{\pcharge{0}{O}}
          (-[:-60]\textcolor{Green}{\ce{CH2CH3}})
          -[:60]\textcolor{Green}{H}
        }
        \arrow{<=>[
          \chemfig{
            @{bromine}\textcolor{Red}{\ce{Br-}}
          }
        ][]}
        \chemfig{
          \ce{CH3CH2}\textcolor{Red}{Br}
        }
        \+
        \chemfig{
          \textcolor{Green}{\ce{HOCH2CH3}}
        }
      }
      \chemmovearrow{
        \reactionfigarrow{oxygen1}{60}{proton}{180}{Red}
        \reactionfigarrow{bond}{-100}{oxygen2}{-100}{Red}
        \reactionfigarrow{bromine}{150}{carbon}{45}{Red}
      }
    }{Ether cleavage with hydrohalic acids}{fig:filler11}{0.8}

  \end{core}

  \begin{core}{\ce{R_{prim}-O-R_{sec}} Reaction}
    
    \ce{R_{prim}-O-R_{sec}} ethers react either through \sno or \snt depending on the nucleophile strength (\textit{under acidic conditions}). When paired with a \textbf{good nucleophile}, the attack occurs on the \ce{R_{prim}} side carbon via \snt to avoid the steric hinderance of \ce{R_{sec}}.

    \reactionfig{
      \scheme{
        \chemfig{
          [:0]
          -[:30]
          (-[:90])
          -[:-30]@{oxygen1}\textcolor{Red}{O}
          -[:30]
          -[:-30]
        }
        \arrow{->[
          \chemfig{
            [:0]@{hydrogen}\textcolor{Blue}{H}
            -[:0,0.7]\textcolor{Red}{I}
          }
        ]}
        \chemname{
          \chemfig{
            [:0,,,,Green]
            -[:30]
            (-[:90])
            -[:-30]@{oxygen2}\textcolor{Green}{\pcharge{90}{O}}
            (-[:-90]\textcolor{Green}{H})
            -[@{bond}:30]@{carbon}
            -[:-30,,,,Black]
          }
        }{\textit{less hindered on right side}}
        \arrow{->[
          \chemfig{
            @{iodine}\textcolor{Red}{\ce{I-}}
          }
        ][]}
        \chemfig{
          [:0]
          -[:60]
          (-[:120])
          -[:0]\textcolor{Red}{O}H
        }
        \+
        \chemfig{
          [:0]\textcolor{Red}{I}
          -[:-60]
          -[:-90]
        }
      }
      \chemmovearrow{
        \reactionfigarrow{oxygen1}{60}{hydrogen}{150}{Red}
        \reactionfigarrow{bond}{-30}{oxygen2}{-45}{Red}
        \reactionfigarrow{iodine}{100}{carbon}{80}{Red}
      }
    }{Mixed ether reaction with good nucleophile (\snt)}{fig:filler12}{0.8}

    When paired with a \textbf{poor nucleophile}, an \sno reaction occurs at the carbocation of \ce{R_{sec}} (\textit{stabalized by hyperconjugation}).

    \myfig{0.65}{ethers_figures/er_2b}{Mixed Ether Reaction with Poor Nucleophile (SN1)}{fig:er_2b}

  \end{core}

  \begin{core}{t-Bu Ether Hydrolysis}

    t-Bu ether hydrolysis is the reverse of t-Bu ether formation covered in Core Concept \ref{core:tbu}. This reverse process is the "deprotection" of alcohol \ce{ROH} by the t-Bu cation via acid. The resulting t-Bu cation, when exposed to heat, undergoes E1 to exaporate into an alkene.

    \myfig{0.65}{ethers_figures/er_3}{Hydrolysis and Deprotection of t-Butyl Ethers}{fig:er_3}
    
  \end{core}

  \begin{core}{Ether Ring Opening}
    
    Ether ring openings release ring strain by a \ce{Nu-} attacking directly one of the two carbons adjacent to the oxygen. This is a regioselective process, meaning \snt occurs at the less hindered relevant carbon.

    \myfig{0.9}{ethers_figures/er_4a}{Regioselective Ether Ring Opening}{fig:er_4a}

    Many \ce{Nu-} will work. The following example shows that ether ring openings are regioselective and also \textbf{stereocontrolled}.

    \myfig{0.4}{ethers_figures/er_4b}{Stereocontrolled Ether Ring Opening}{fig:er_4b}

    With \textbf{neutral nucleophiles}, acidic conditions are required to activate the ether toward nucleophilic attack. Without \ce{H+}, the reaction will not occur.

    \myfig{0.45}{ethers_figures/er_4c}{Acid-Activated Ether Ring Opening with Neutral Nucleophile}{fig:er_4c}
    
    One important distinction in the neutral nucleophile mechanism is regioselectivity towards the more hindered carbon. In asymmetrical ethers, the nucleophile preferentially attacks the more substituted carbon, as it better stabilizes the partial positive charge (δ$^+$) that develops through hyperconjugation and inductive effects---causing the C-O bond on that side to stretch and weaken. This gives the mechanism a hybrid \sno/\snt character: it resembles \snt in that nucleophile attack and leaving group departure occur simultaneously, yet resembles \sno in that the partial positive charge mimics a carbocation.

    \myfig{0.55}{ethers_figures/er_4d}{Regioselectivity in Acid-Catalyzed Ring Opening of Asymmetrical Ethers}{fig:er_4d}

  \end{core}

\section{Sulfur Analogs: \textit{A Brief Lesson}}

  \begin{definition}{Sulfur Analogs}
    
    Sulfur is directly below oxygen in the periodic table, and thus behaves very similarly.

    Thiols (\ce{RSH}) are \textbf{more acidic than alcohols} because the \ce{S-H} bond is weaker due to sulfur's higher-energy molecular orbitals. The lower electronegativity difference in \ce{S-H} compared to \ce{O-H} also results in \textbf{weaker hydrogen bonding}, making thiols less polar overall. \textit{This is illustrated by \ce{H2S}, which is a gas at 25$^\circ$C with a boiling point of $-60^\circ$C, in stark contrast to \ce{H2O}.} 

    Despite being less basic than \ce{RO-}, \ce{RS-} is a stronger nucleophile due to its greater \textbf{polarizability}---note that polarizability is distinct from polarity. Polarizability describes how easily an electron cloud can be distorted by an external force (covered in EXTERNAL LABEL PLACEHOLDER TO NUC SUB AND ELIM FILE). Larger atoms like sulfur are more polarizable than oxygen because their diffuse, loosely-held electron clouds are more easily deformed. This allows \ce{RS-} to stretch and distort toward an electrophile, enhancing its nucleophilicity despite its lower basicity.

  \end{definition} 

  \begin{core}{Sulfur Reactivity}
    
    The better nucleophility of \ce{RS-} eliminates the issue of E2 reactions when paired with \ce{R_{sec}X} like \ce{RO-} has.

    \myfig{0.6}{ethers_figures/sulfur_e2_elimin}{sulfer e2 elimination}{fig:sulfere2_elim}
    
    Even neutral \ce{RSR{'}} readily undergoes \snt, whereas dimethyl ether (\ce{CH3OCH3}) does not react under the same conditions. When \ce{RSR{'}} attacks an \ce{RX}, the resulting sulfonium ion (\ce{R3S+}) becomes susceptible to attack by another nucleophile — the sulfide, once a good nucleophile, has now become a good leaving group, analogous to \ce{H2O}. This makes \textbf{sulfonium salts} (\ce{R3S+}) effective \textbf{alkylating agents}, meaning they can transfer an alkyl group to an incoming nucleophile.

    \myfig{0.75}{ethers_figures/sulfur_goodnul}{sulfer good nul}{fig:sulfer_goodnul}

  \end{core}

\end{document}
